\subsection{Bối cảnh nghiên cứu}

Sự phát triển vượt bậc của các Mô hình Ngôn ngữ Lớn (Large Language Models - LLMs) đã thúc đẩy một cuộc cách mạng trong lĩnh vực Xử lý Ngôn ngữ Tự nhiên (NLP), đặc biệt là trong các hệ thống đối thoại tương tác (Conversational Agents). Để vượt qua giới hạn về tri thức tĩnh của LLM, kỹ thuật Sinh văn bản tăng cường truy xuất (Retrieval-Augmented Generation - RAG) đã trở thành tiêu chuẩn công nghiệp, cho phép mô hình trích xuất thông tin từ kho tài liệu ngoại vi trước khi tổng hợp câu trả lời.

Tuy nhiên, khi triển khai vào các hệ thống doanh nghiệp lõi (Enterprise Tech Core) đòi hỏi tính nguyên tử và khắt khe về mặt logic như Tư vấn, Đặt chỗ và Lập lịch trình du lịch, kiến trúc RAG nguyên bản (Native RAG) đã bộc lộ những khiếm khuyết cơ bản về mặt lý thuyết thông tin và cấu trúc suy luận. Việc thiếu vắng một cơ chế tương tác với dữ liệu thay đổi theo thời gian thực (Real-time Dynamic Data) và khả năng xử lý các ràng buộc chéo (Cross-constraints) đã hạn chế nghiêm trọng khả năng ứng dụng của LLM, biến chúng từ một trợ lý thông minh thành một hệ thống tiềm ẩn rủi ro sai lệch thông tin nghiêm trọng.
